\documentclass{article}
\usepackage{amsmath, amssymb, amsthm}
\usepackage{hyperref}
\usepackage{geometry}
\geometry{margin=1in}

\title{Erd\H{o}s Problem \#177}
\date{Last updated: 2025-08-31}
\author{Source: erdosproblems.com}

\begin{document}
\maketitle

\noindent\textbf{Status:} OPEN \quad \textbf{No prize}

\noindent\textbf{Tags:} discrepancy, arithmetic progressions

\medskip

\noindent\textbf{Problem Statement:}

Find the smallest $h(d)$ such that the following holds. There exists a function $f:\mathbb{N}\to\{-1,1\}$ such that, for every $d\geq 1$,\[\max_{P_d}\left\lvert \sum_{n\in P_d}f(n)\right\rvert\leq h(d),\]where $P_d$ ranges over all finite arithmetic progressions with common difference $d$.




Cantor, Erd\H{o}s, Schreiber, and Straus \cite{Er66} proved that $h(d)\ll d!$ is possible. Van der Waerden's theorem implies that $h(d)\to \infty$. Beck \cite{Be17} has shown that $h(d) \leq d^{8+\epsilon}$ is possible for every $\epsilon&#62;0$. Roth's famous discrepancy lower bound \cite{Ro64} implies that $h(d)\gg d^{1/2}$.




References

[Be17] Beck, J\'{o}zsef, A discrepancy problem: balancing infinite dimensional vectors . Number theory-Diophantine problems, uniform distribution and applications (2017), 61-82.

[Er66] Erd\H{o}s, P\'{a}l, Remarks on number theory. {V}. {E}xtremal problems in number theory. {II} . Mat. Lapok (1966), 135--155.

[Ro64] Roth, K. F., Remark concerning integer sequences . Acta Arith. (1964), 257-260.


\bigskip
\noindent\small{Source: \url{https://www.erdosproblems.com/177}}
\end{document}
